% a mashup of hipstercv, friggeri and twenty cv
% https://www.latextemplates.com/template/twenty-seconds-resumecv
% https://www.latextemplates.com/template/friggeri-resume-cv

\documentclass[lighthipster]{simplehipstercv}
% available options are: darkhipster, lighthipster, pastel, allblack, grey, verylight, withoutsidebar
% withoutsidebar
\usepackage[utf8]{inputenc}
\usepackage[default]{raleway}
\usepackage[margin=1cm, a4paper]{geometry}


\usepackage{afterpage}
\usepackage{atbegshi}

\hypersetup{colorlinks=true, urlcolor=pblue, linkcolor=black}

\definecolor{sidebargrey}{rgb}{0.9,0.9,0.9}


%------------------------------------------------------------------ Variablen

\newlength{\rightcolwidth}
\newlength{\leftcolwidth}
\setlength{\leftcolwidth}{0.23\textwidth}
\setlength{\rightcolwidth}{0.75\textwidth}

%------------------------------------------------------------------
%\title{New Simple CV}
%\author{\LaTeX{} Ninja}
%\date{June 2019}

\pagestyle{empty}
\begin{document}


\thispagestyle{empty}
%-------------------------------------------------------------

% light grey bar on the left edge of every page, full height; on page 1 the
% dark header is drawn on top of it and covers its upper part
\AtBeginShipout{\AtBeginShipoutUpperLeft{\color{sidebargrey}\rule[-\paperheight]{2.1cm}{\paperheight}}}

\simpleheader{headercolour}{David}{Hägele}{PhD Student \hspace{1em} \icon{\faEnvelopeO}{white}{} \href{mailto:haegele@mpia.de}{haegele@mpia.de}}{white}

%------------------------------------------------

% this has to be here so the paracols starts..
\subsection*{}
\vspace{4em}

\setlength{\columnsep}{1.5cm}
\columnratio{0.02}[0.98]
\begin{paracol}{2}
\hbadness5000

\switchcolumn

\small
\section*{Education}

\begin{tabular}{r| p{0.67\textwidth} c}
    \cvevent{Oct 2025--present}{Doctor of Philosophy in Astronomy}{Max Planck Institute for Astronomy}{Heidelberg, Germany}{Describe it here!}{} \\
    \cvevent{Aug 2023--Sep 2025}{Master of Science in Physics}{Heidelberg University, Max Planck Institute for Astronomy}{Heidelberg, Germany}{Focus on exoplanets, astronomy, and environmental physics. Thesis with Paul Mollière on a Python package to calculate high-resolution opacities.}{}\\
    \cvevent{Aug 2023--Jan 2024}{Erasmus Semester}{Uppsala University}{Uppsala, Sweden}{Lectures on planetary and space physics and astronomy. Research project with Nikolai Piskunov on principal component analysis-based removal of telluric lines from generated synthetic transmission
    spectra of gas giant exoplanets.}{}\\
    \cvevent{Sep 2020--Aug 2023}{Bachelor of Science in Physics}{Heidelberg University, Max Planck Institute for Astronomy }{Heidelberg, Germany}{Focus on astronomy and exoplanets. Thesis with Maria Steinrück on an enhanced parameterization of 1D vertical mixing in hot Jupiter atmospheres.}{}\\
\end{tabular}


\vspace{0.4em}


\section*{Publications}
\textbf{As first-author}
\vspace{0.5em}
\begin{enumerate}
    \item Haegele, D. B., Mollière, P., Min, M. (2026): line racer: rapid calculation of exoplanet radiative opacities. Submitted to the Journal of Open Source Software. 
\end{enumerate}

\vspace{0.4em}

\textbf{As co-author}
\vspace{0.5em}
\begin{enumerate}
    \item Steinrueck, M. E. et al. (incl. \textbf{David Haegele}, submitted): Limb asymmetries on WASP-39b: A Multi-GCM Comparison of Chemistry, Clouds, and Hazes. Submitted to The Astrophysical Journal, arXiv:2509.21588.
    \item Espinoza, N., Steinrueck, M.E., Kirk, J. et al. (incl. \textbf{David Haegele}) Inhomogeneous terminators on the exoplanet WASP-39 b. Nature \textbf{632}, 1017–1020 (2024). \href{https://doi.org/10.1038/s41586-024-07768-4}{https://doi.org/10.1038/s41586-024-07768-4}
\end{enumerate}


\vspace{0.4em}


\section*{Software}
\begin{tabular}{r| p{0.67\textwidth} c}
    \cvevent{line racer}{Main developer}{Open source Python package}{\href{https://line-racer.readthedocs.io/en/latest/}{line-racer.readthedocs.io}}{Designed and wrote the package from scratch: rapid calculation of high-resolution radiative opacities for exoplanet atmospheres (JOSS paper submitted).}{} \\
    \cvevent{petitRADTRANS}{Maintainer \& core developer}{Open source Python package}{\href{https://petitradtrans.readthedocs.io/en/latest/}{petitradtrans.readthedocs.io}}{Part of the active development team of the widely used package for calculating exoplanet spectra and running atmospheric retrievals; maintain the code base and develop new features.}{} \\
\end{tabular}


\vspace{0.4em}


\section*{Presentations}

\begin{itemize}
    \item Spectroscopy of Exoplanets Over All Wavelength, contributed talk, June 2025
    \item A New Era in Exoplanet Atmosphere Observation and Characterisation (summerschool), poster presentation, July 2025
\end{itemize}


\vspace{0.4em}


\section*{Teaching experience}

\begin{itemize}
    \item Abitur preparation tutor for math and physics (2020, 2022, 2024, 2025): Repetition of the content of two years and exercises in preparation for the final exams. Groups of 12 students for one week of 6 hours per day.
\end{itemize}

\vfill{}
\vfill{} % Whitespace before final footer

\end{paracol}

\end{document}
